\section{Acquisition, repeated cells, preparation and the full service bill}
\label{sec:acquisition}
The decision uses coarse certificates, not a finely calibrated model. Their acquisition remains part of an end-to-end claim. Already valid engineering or previously acquired certificates can be reused, but their provenance and costs are not erased.

\subsection{Coarse probabilities and actual operating clocks from paired holds}
At one actual menu environment and neighbor context, let $r$ and $q$ be the unknown two-state clock and equilibrium failure probability. Pin to zero and read after a hold to obtain a $0\to1$ indicator; separately pin to one for a $1\to0$ indicator. At a common exact hold $t$ their means are
\begin{equation}
 a=q(1-e^{-rt}),\qquad b=(1-q)(1-e^{-rt}).
 \label{eq:calpair}
\end{equation}
Neither experiment waits for a locally rare natural starting bit. The two paid pins provide it. Their sum exposes the refresh clock without a Boltzmann suppression, and their ratio exposes $q$.

The hold lengths need not be identical. Suppose each is exogenous to its current focal path and belongs to $[t_-,t_+]$, with $t_->0$. They may depend on previous experimental history. Let simultaneous confidence intervals for the conditional average indicator means be $[a_-,a_+]$ and $[b_-,b_+]$, including all pin/read/context mean biases. Put $c_t=t_+/t_-$. Concavity of $1-e^{-rt}$ implies that the ratio of any two admissible hold responses is between $c_t^{-1}$ and $c_t$. Thus valid intervals include
\begin{align}
 q&\in\left[\frac{a_-}{a_-+c_t b_+},\,
             \frac{c_t a_+}{c_t a_++b_-}\right],\label{eq:qinterval}\\
 r&\in\left[\frac{-\log(1-d_-^{\rm obs})}{t_+},\,
             \frac{-\log(1-d_+^{\rm obs})}{t_-}\right],
 \quad d_-^{\rm obs}=\max(0,a_-+b_-),\quad
 d_+^{\rm obs}=\min(1,a_++b_+).
 \label{eq:rinterval}
\end{align}
A zero denominator gives the conservative endpoint zero or one; $d_+^{\rm obs}=1$ gives an infinite upper clock bound. A supplied physical upper rate bound remains useful. No finite observation can certify an arbitrary unbounded rate above a saturated detector's range.

For $m$ observations of either starting bit, martingale Hoeffding gives statistical radius $\sqrt{\log(2C/\delta_{
\rm cal})/(2m)}$ around its conditional average; add the certified conditional mean bias. This is valid for reused cells with the declared memoryless conditional response. It does not infer memorylessness from a short record.

Use one menu item's probability intervals to certify the common equilibrium law. Enumerate the declared local contexts. For neighboring contexts $s,s'$ differing at $j$, form the maximum difference allowed by their $q$ intervals. Maximize over the remaining local bits to obtain $\overline c_{ij}$. Verify
\begin{equation}
 \max_{i,s}q_i^+(s)\le\bar q<1/2,\qquad
 \max_i\sum_j\overline c_{ij}\le1-\alpha.
 \label{eq:coarsecertificate}
\end{equation}
Clock lower bounds must hold in every declared context of every menu item. Their enumeration costs $C_{
\rm ctx}=K\sum_i2^{\deg(i)}$, not its logarithm. Common equilibrium across items is a structural premise, not a consequence of similar fitted probability intervals. These bounds give $\gamma=\alpha r_-$.

If the certified probability intervals also give $q_i(s)\ge q_->0$, then every configuration has $\pi(x)\ge q_-^N$, since $1-q_i(s)\ge1-\bar q>q_-$ after reducing $q_-$ if necessary. This supplies the mass floor for relative preparation. Otherwise the accepted bounded-energy learner can supply that floor at its separately paid field/dwell cost. The decision itself uses no lower bound on $q_i$.

Under fixed strict margins, all intervals above can be obtained with constant per-context probability accuracy. If $\Delta_{
\rm cal}$ denotes a sufficient additive probability accuracy, the coarse acquisition uses
\begin{equation}
 O\left(C_{\rm ctx}\Delta_{\rm cal}^{-2}
            \log\frac{C_{\rm ctx}}{\delta_{\rm cal}}\right)
 \label{eq:calcost}
\end{equation}
pins and binary reads, and at most $t_+$ times that count in free-cell exposure. The actual hold choice and all context, timer and field costs are added. A sequence of decreasing accuracies with summable confidence can discover strict slack; an exact boundary need not terminate. A promised positive slack gives a finite stated schedule. In particular, no global $1/p$ factor is introduced, but rare local responses, large degree or small margins can make calibration expensive.

\subsection{What a local service costs}
\label{sec:services}
For the preferred boundary branch, a single call costs
\begin{equation}
 C_{\rm setup}(d,u)+C_{\rm pin}(\beta)+C_{\rm field}(\xi_u)
 +C_{\rm slew}+h(1+\omega)+C_{\rm read}(\beta)+C_{\rm timer}(h\omega)
 +C_{\rm bits}.
 \label{eq:physicalservice}
\end{equation}
These symbols denote actual validated services, not unit-cost promises. Context pins are paid, including setting the neighbors to the failed value. The system itself is never prepared globally failed: only one focal local copy and its declared neighborhood are set. The operating environment is reproduced at every experiment and on deployment. If a small field error can change an unknown clock discontinuously, field precision alone gives no clock guarantee; a reproduced environment or a proved operating-rate envelope is required.

Boundary trials need only conditional probability bias bounded by the decision-level allowance in Section~\ref{sec:boundary}. It does not shrink as the number of trials grows. Dependence from repeated cells is handled by the conditional concentration theorem. The exogenous timer may drift within its declared range as a function of past observations. A separate weaker bound covers a stopping time responsive to the current focal path. Neither permits unbounded detector delay or missed exits without its stated allowance.

The return branch requires a fixed reversible ideal $P_h$. A finite cap bounds its total calls by $D_{
\rm calls}=KnJ$. Conditional pin, read, numerical and selector errors with sum at most $\delta_{
\rm service}$ couple it to that ideal experiment except with that probability. This branch therefore may need finer service than the direct boundary branch. The distinction is intentional. Finite random-coordinate selection is implemented by binary rejection with a cap; its bias can be charged at the mean level for the boundary branch, or in the total coupling ledger for returns. Integer counters, rational interval arithmetic and output confidence bits are paid.

There is no universal claim that $C_{\rm pin}(\beta)$ or $C_{\rm read}(\beta)$ grows logarithmically. If a given instrument supports a stable captured bit and conditionally bounded independent read faults below one half, repeated reads and majority give that familiar logarithmic amplification. The capture and fault assumptions are additional supplied hardware in that realization. The theorem otherwise keeps the actual service function.

\subsection{The separate stronger relative-equilibrium initialization}
The mean-only boundary decision can instead use Section~\ref{sec:anystart}.
This subsection retains the full relative-law preparation for the return
fallback and the older forecasts that genuinely require it.
For a selected menu item let the repairable operating system mix, inspect it, retain a safe state and otherwise retry. The process must remain available through failed states during preparation; this is not free preparation of an irreversibly destroyed specimen. The stationary safe mass is at least $1-\bar p\ge1/2$.

With $\pi_{
\min}\ge q_-^N$, reversible spectral mixing gives, from any initial state,
\[
 \left|\frac{P_t(x,y)}{\pi(y)}-1\right|\le e^{-\gamma t}/\pi_{
\min}.
\]
Choose $t_{
\rm prep}\ge\gamma^{-1}\log[1/(\xi\pi_{
\min})]$ and set $\xi\le\tanh(\beta/2)$. Conditioning a correctly inspected safe state then yields
\begin{equation}
 e^{-\beta}\le\frac{d\nu_u}{d\pi(\cdot\mid A)}\le e^\beta.
 \label{eq:prep}
\end{equation}
A finite $O(\log(1/\delta_{
\rm prep}))$ retry cap suffices at the stated safe mass, with explicit preparation abort. Inspection and hold errors must be included using the accepted finite relative-preparation construction. A smaller safe-mass floor would require its inverse in the retry cost. Here the locally certified $\bar q$ supplies the safe-mass slack directly.

This pointwise density bound, not just TV, is required for an unbounded mean:
$e^{-\beta}M_u\le\E_{
\nu_u}\tau_b\le e^\beta M_u$.
If the actual deployed generator additionally has a common-invariant form discrepancy $e^{\pm\zeta}$ from the calibrated one, the ratio guarantee loses at most $e^{-2(\beta+\zeta)}$. A mean-zero preparation-abort outcome multiplies an unconditional lower guarantee by its success probability; for a conditional-on-success mission, the abort probability is reported separately. General changes to the invariant law do not enjoy this cancellation and require the earlier acquired uncertainty machinery.

Only the selected action is physically prepared. Acquisition precedes the mission. Protecting a live installation during acquisition needs its own priced hold, shutdown or exposure analysis. The intervention is not instantaneous feedback: computation, installation, field travel, finite slew, settling and inspection all take time. The deployment bill is
\begin{equation}
 C_{
\rm install}+C_{
\rm ramp}+C_{
\rm settle}
 +O\left(\log(1/\delta_{
\rm prep})\,[t_{
\rm prep}+C_{
\rm inspect}(N)]\right).
 \label{eq:deployment}
\end{equation}
Under this displayed sufficient relative-mixing bound, the free-evolution wall time is $\widetilde O(N)$ at fixed margins. It acts on an entire $N$-component installation and therefore uses $\widetilde O(N^2)$ summed component exposure. The inspection and other service functions are additional. It is not a lower bound on possible preparation time. The acquired-information separation does not justify silently dropping this deployment cost.

\subsection{Resource scaling and nonclaims}
At fixed degree, clock contrast, reliability margin, decision accuracy and confidence, coarse calibration has $\widetilde O(KN2^d)$ local calls. The boundary decision has a number of local calls independent of $N$; the return branch has $O(K\operatorname{polylog}N)$, with the explicit accuracy and numerical margin factors in its schedule. A sparse set of recovered indices represents a return configuration; initializing $b$ requires an empty set, not $N$ physical failures. The known interaction graph and full menu description still occupy their actual input size. Loading them can cost $\Omega(N)$ even when a subsequent decision samples only a few coordinates.

These are local experiment and bit-work statements, not claims of constant wall time for a growing apparatus. Parallel acquisition pays the number of cells, controllers and detectors; the prediction lower bound charges their summed exposure. Large sufficient calibration constants can dominate the short decision. The numerical executions below do not establish economical general hardware.
